\documentclass[12pt,a4paper]{article}
\usepackage{amsmath,amssymb}
\usepackage{graphicx}
\usepackage{hyperref}
\usepackage{booktabs}
\usepackage{geometry}
\geometry{margin=1in}

\title{Lensing-Dynamics Mass Discrepancy as a Redshift-Dependent\\
Signature of Informational Gravity}

\author{Heath W.\ Mahaffey\\
\small Independent Researcher, Entiat, WA 98822, USA\\
\small \texttt{hmahaffeyges@gmail.com}}

\date{February 25, 2026\\
\small \url{https://github.com/hmahaffeyges/IAM-Validation}\\
\small Zenodo: \href{https://doi.org/10.5281/zenodo.18750699}{10.5281/zenodo.18750699}}

\begin{document}

\maketitle


\begin{abstract}
We analyze gravitational lensing in the dual-sector realization of the
Informational Actualization Model (IAM), in which gravitational decoherence
renormalizes the effective expansion rate governing matter perturbations
while leaving the background Friedmann evolution and photon sector unchanged.
In the standard $\mu$--$\Sigma$ parameterization of modified gravity,
IAM predicts $\mu(a) < 1$ and $\Sigma(a) = 1$ exactly at linear order.

We show explicitly that the linearized Einstein equations for the metric
potentials remain unmodified and that the informational field does not
introduce anisotropic stress or perturbative degrees of freedom.
Consequently, weak lensing observables are affected only indirectly through
the altered growth history of $\delta_m$, while the relation between the
Weyl potential and matter density is preserved.

We discuss the implications for $S_8$ measurements and current weak lensing
constraints, emphasizing that IAM produces a modest suppression of growth
consistent in direction with existing survey trends, without introducing
additional free parameters beyond those of $\Lambda$CDM.
\end{abstract}


%-------------------------------------------------------------------
\section{Introduction}
%-------------------------------------------------------------------

\sloppy
Galaxy cluster mass estimates obtained through gravitational lensing consistently
exceed those obtained through dynamical methods --- galaxy velocity dispersions
and X-ray hydrostatic equilibrium --- by 10--30\%~\citep{Mahdavi2013,Israel2014,
Hoekstra2015,vonDerLinden2014,Smith2016}. This discrepancy is conventionally
attributed to hydrostatic bias: the assumption that the intracluster medium is
in perfect hydrostatic equilibrium, neglecting non-thermal pressure support from
bulk motions, turbulence, and cosmic rays. The standard parametrization is:
\begin{equation}
M_\mathrm{true} = \frac{M_{X\text{-ray}}}{1-b}
\end{equation}
with observational estimates ranging from $b \approx 0.1$ to $b \approx 0.4$
depending on methodology~\citep{Nagai2007,Rasia2012,Biffi2016}.

The hydrostatic bias explanation has several unsatisfying features. The required
bias is larger than predicted by hydrodynamic simulations ($b \approx 0.1$--$0.15$
from~\citep{Lau2009,Nelson2014}). The bias shows no clear dependence on cluster
properties that would be expected if it were purely a gas physics effect. And
the tension between Planck~SZ cluster counts and Planck~CMB cosmology --- which
requires $b \approx 0.4$, far above simulation predictions --- suggests the
explanation may not be purely astrophysical~\citep{PlanckClusters2016}.

In this paper, we show that the Informational Actualization
Model~\citep{Mahaffey2026a,Mahaffey2026b} naturally produces a lensing-dynamics
mass discrepancy with the correct magnitude and a distinctive redshift dependence
that no other proposed explanation shares. The prediction follows directly from
IAM's unique signature: $\mu < 1$ (weakened gravitational coupling for matter)
with $\Sigma = 1$ (unmodified light deflection). Full Boltzmann validation
through 15 converged Planck~2018 MCMC chains confirms $\Delta\chi^2 = +0.54$
relative to $\Lambda$CDM with $\sigma_8 = 0.800$~\citep{Mahaffey2026a}.

%-------------------------------------------------------------------
\section{The IAM Prediction}
%-------------------------------------------------------------------

\subsection{The $\mu$--$\Sigma$ framework}

In the standard modified gravity
parametrization~\citep{Pogosian2016,Amendola2008}, deviations from General
Relativity are captured by two functions of scale factor:
\begin{align}
\nabla^2\Psi &= -4\pi G a^2\, \mu(a)\, \rho\delta
\quad \text{(Poisson equation, dynamics)} \\
\nabla^2(\Psi + \Phi) &= -8\pi G a^2\, \Sigma(a)\, \rho\delta
\quad \text{(lensing potential)}
\end{align}
where $\Psi$ is the Newtonian potential governing matter dynamics and
$\Phi + \Psi$ is the lensing potential governing photon deflection.
In GR, $\mu = \Sigma = 1$.

\subsection{IAM values}

IAM predicts~\citep{Mahaffey2026a,Mahaffey2026b}:
\begin{equation}
\mu(a) = \frac{H^2_\Lambda(a)}{H^2_\Lambda(a) + \beta_m\, E(a)}, \qquad
\Sigma(a) = 1 \quad \text{(exact, all redshifts)}
\label{eq:mu}
\end{equation}
where $H^2_\Lambda(a) = \Omega_m a^{-3} + \Omega_\Lambda$,
$\beta_m = \Omega_m/2 = 0.1575$ (virial theorem, derived),
and $E(a) = \exp(1-1/a)$ (activation function, derived from horizon
thermodynamics). No parameters are fitted.

The physical origin of $\Sigma = 1$ is fundamental: photons travel on null
geodesics ($d\tau = 0$) and do not experience gravitational decoherence. The
informational entropy contribution that modifies gravity for matter is invisible
to light. Lensing sees the full, unmodified gravitational potential.

\subsection{Mass ratio prediction}

A cluster's dynamical mass is inferred from matter responding to the $\Psi$
potential:
\begin{equation}
M_\mathrm{dyn} \propto \Psi \propto \mu \times M_\mathrm{true}
\end{equation}
A cluster's lensing mass is inferred from photons responding to the
$(\Psi + \Phi)/2$ potential:
\begin{equation}
M_\mathrm{lens} \propto \frac{\Psi+\Phi}{2} \propto \Sigma \times M_\mathrm{true}
= M_\mathrm{true}
\end{equation}
Therefore:
\begin{equation}
\boxed{\frac{M_\mathrm{lens}}{M_\mathrm{dyn}} = \frac{\Sigma}{\mu} = \frac{1}{\mu(z)}}
\label{eq:ratio}
\end{equation}
This is a zero-parameter prediction. The entire $M_\mathrm{lens}/M_\mathrm{dyn}(z)$
curve is determined by $\beta_m = \Omega_m/2$ and $E(a) = \exp(1-1/a)$.

\begin{table}[h]
\centering
\caption{Predicted lensing-to-dynamical mass ratio $M_\mathrm{lens}/M_\mathrm{dyn}
= 1/\mu(z)$ as a function of redshift. Computed from Equation~\ref{eq:mu} with
$\beta_m = 0.1575$ and Planck~2018 parameters. The $15.7\%$ excess at $z=0$
declines monotonically to zero at high redshift.}
\label{tab:ratio}
\begin{tabular}{cccc}
\toprule
Redshift $z$ & $\mu(z)$ & $M_\mathrm{lens}/M_\mathrm{dyn}$ & Excess (\%) \\
\midrule
0.0 & 0.864 & 1.157 & 15.7 \\
0.1 & 0.886 & 1.129 & 12.9 \\
0.2 & 0.905 & 1.105 & 10.5 \\
0.3 & 0.922 & 1.085 &  8.5 \\
0.5 & 0.948 & 1.055 &  5.5 \\
0.7 & 0.966 & 1.035 &  3.5 \\
1.0 & 0.982 & 1.018 &  1.8 \\
1.5 & 0.994 & 1.006 &  0.6 \\
2.0 & 0.998 & 1.002 &  0.2 \\
3.0 & 1.000 & 1.000 &  0.0 \\
\bottomrule
\end{tabular}
\end{table}

\subsection{Explicit functional form}

Substituting Equation~\ref{eq:mu} into Equation~\ref{eq:ratio}:
\begin{equation}
\frac{M_\mathrm{lens}}{M_\mathrm{dyn}}(a)
= 1 + \frac{\beta_m\, E(a)}{H^2_\Lambda(a)}
= 1 + \frac{\beta_m\, \exp(1-1/a)}{\Omega_m a^{-3} + \Omega_\Lambda}
\label{eq:ratio_explicit}
\end{equation}
This is the observable prediction: a monotonically declining function of $a$
(increasing function of $z$), fully determined by Planck~2018 parameters and the
virial theorem, with no fitted parameters.

%-------------------------------------------------------------------
\section{Comparison with Alternative Explanations}
%-------------------------------------------------------------------

\subsection{Hydrostatic bias}

The standard hydrostatic bias $b$ produces a constant mass ratio
$M_\mathrm{lens}/M_\mathrm{dyn} = 1/(1-b) \approx 1.2$ for $b = 0.17$,
independent of redshift by construction. Some redshift dependence might arise
from cluster evolution (relaxation state), but this would correlate with
dynamical state indicators and would not follow the specific functional form
of Equation~\ref{eq:ratio_explicit}.

\textit{Distinguishing test:} Plot $M_\mathrm{lens}/M_\mathrm{dyn}$ vs.\ $z$.
Hydrostatic bias predicts a flat line or scatter with no systematic trend. IAM
predicts a specific declining curve.

\subsection{MOND and MONDian theories}

Modified Newtonian Dynamics and its relativistic extensions predict enhanced
gravitational effects at low accelerations ($a < a_0 \approx 1.2\times
10^{-10}$~m\,s$^{-2}$). The mass discrepancy depends on the acceleration profile
of each individual cluster, not on cosmological redshift.

\textit{Distinguishing test:} At fixed cluster mass and morphology, MOND predicts
$M_\mathrm{lens}/M_\mathrm{dyn}$ depends on position within the cluster. IAM
predicts it depends only on the cluster's observed redshift, independent of
internal properties.

\subsection{Standard modified gravity ($f(R)$, nDGP)}

Most modified gravity theories predict $\mu > 1$ --- gravity is stronger than
GR, not weaker. This produces $M_\mathrm{dyn} > M_\mathrm{lens}$, the opposite
of what is observed. The few models that allow $\mu < 1$ typically also predict
$\Sigma \neq 1$, which conflicts with the unmodified CMB lensing power spectrum
measured by Planck~\citep{Planck2020lens}.

IAM is unique in the current literature in predicting $\mu < 1$ with $\Sigma = 1$
exactly, arising from a physical mechanism (decoherence asymmetry between
timelike and null worldlines) rather than parameter tuning.

\subsection{Summary of distinguishing signatures}

\begin{table}[h]
\centering
\small
\caption{Predicted lensing-dynamics signatures for competing models.
IAM is the only model predicting $\mu < 1$ with $\Sigma = 1$ exactly and
a specific, parameter-free redshift dependence.}
\label{tab:models}
\begin{tabular}{p{3.0cm}p{2.8cm}p{2.8cm}p{2.8cm}}
\toprule
Model & $M_\mathrm{lens}/M_\mathrm{dyn}$ at $z=0$ & $z$-dependence & $\mu$ vs.\ $\Sigma$ \\
\midrule
GR + hydrostatic bias & $\sim 1.2$ (fitted $b$) & Flat or scatter & $\mu = \Sigma = 1$ \\[4pt]
MOND & Varies with $a/a_0$ & Acceleration-dependent & $\mu > 1$, $\Sigma > 1$ \\[4pt]
$f(R)$ gravity & $< 1$ ($M_\mathrm{dyn} > M_\mathrm{lens}$) & Scale-dependent & $\mu > 1$, $\Sigma > 1$ \\[4pt]
nDGP gravity & $< 1$ & $z$-dependent & $\mu > 1$, $\Sigma \geq 1$ \\[4pt]
\textbf{IAM} & \textbf{1.157} & \textbf{Eq.~\ref{eq:ratio_explicit}, zero params} & $\mu < 1$, $\Sigma = 1$ \\
\bottomrule
\end{tabular}
\end{table}

%-------------------------------------------------------------------
\section{Existing Observational Evidence}
%-------------------------------------------------------------------

\subsection{Planck SZ cluster counts}

The Planck~SZ cluster count analysis~\citep{PlanckClusters2016} found a
significant tension between the number of detected clusters and the prediction
from Planck~CMB cosmology. Resolution requires either $S_8 \approx 0.78$ (well
below the CMB value of $0.83$) or a hydrostatic bias $b \approx 0.4$ (far above
simulation predictions of $b \approx 0.1$--$0.15$). IAM resolves this naturally:
$\mu < 1$ suppresses $\sigma_8$ from $0.811$ to $0.800$, reducing the predicted
cluster count without requiring extreme hydrostatic bias. The residual discrepancy
is consistent with a moderate $b \approx 0.1$--$0.15$, in line with simulations.

\subsection{Canadian Cluster Comparison Project (CCCP)}

Hoekstra et al.~\citep{Hoekstra2015} compared weak lensing masses to X-ray
hydrostatic masses for 50 massive clusters at $0.15 < z < 0.55$, finding
$M_\mathrm{lens}/M_X = 1.20 \pm 0.12$ on average. The IAM prediction at the
median redshift $z \approx 0.3$ is $M_\mathrm{lens}/M_\mathrm{dyn} = 1.085$,
within the measurement uncertainty. CCCP did not bin by redshift with sufficient
resolution to test the $z$-dependence.

\subsection{Weighing the Giants (WtG)}

Von der Linden et al.~\citep{vonDerLinden2014} performed weak lensing calibration
of Planck cluster masses, finding $M_\mathrm{lens}/M_\mathrm{Planck} = 1.31
\pm 0.11$ for the full sample. The IAM prediction at the sample median
$z \approx 0.25$ is $1.095$, below the WtG central value. However, the WtG
comparison includes Planck mass estimation systematics beyond hydrostatic bias,
making a direct comparison model-dependent.

\subsection{The discrepancy pattern}

Across multiple analyses, the consistent finding is that lensing masses exceed
dynamical or X-ray masses by 10--30\%. The magnitude is consistent with IAM's
$z = 0$ prediction of $15.7\%$. The scatter is consistent with varying cluster
redshifts, dynamical states, and measurement systematics overlaid on a systematic
IAM signal. No existing analysis has performed a high-precision test of the
redshift evolution in the mass discrepancy using a uniform methodology across a
wide redshift baseline. This is the critical test.

%-------------------------------------------------------------------
\section{Predictions for Forthcoming Surveys}
%-------------------------------------------------------------------

\subsection{Euclid}

The Euclid mission~\citep{Euclid2022} will provide weak lensing mass estimates
for $\sim 10^5$ galaxy clusters out to $z \approx 2$. Combined with spectroscopic
velocity dispersions, this enables a direct $M_\mathrm{lens}/M_\mathrm{dyn}$
comparison in fine redshift bins. The IAM predictions from
Equation~\ref{eq:ratio_explicit}:

\begin{table}[h]
\centering
\caption{IAM predictions for $M_\mathrm{lens}/M_\mathrm{dyn}$ in Euclid redshift
bins, computed from Equation~\ref{eq:ratio_explicit} with no free parameters.
The declining trend from $10.5\%$ to $<1\%$ over $0.2 < z < 2.0$ is the
falsifiable signature.}
\label{tab:euclid}
\begin{tabular}{ccc}
\toprule
Redshift bin & $M_\mathrm{lens}/M_\mathrm{dyn}$ (IAM) & Excess (\%) \\
\midrule
$z = 0.2$ & 1.105 & 10.5 \\
$z = 0.5$ & 1.055 &  5.5 \\
$z = 1.0$ & 1.018 &  1.8 \\
$z = 1.5$ & 1.006 &  0.6 \\
$z = 2.0$ & 1.002 &  0.2 \\
\bottomrule
\end{tabular}
\end{table}

If the discrepancy is flat (hydrostatic bias) or acceleration-dependent (MOND),
IAM is challenged. If it follows the declining curve of
Equation~\ref{eq:ratio_explicit}, that is a direct measurement of $\mu(z)$ from
cluster data.

\subsection{Rubin Observatory LSST}

The Legacy Survey of Space and Time~\citep{LSST2009} will provide photometric
cluster catalogs and weak lensing mass maps. The statistical power of
$\sim 200{,}000$ clusters will enable percent-level mass ratio measurements in
narrow redshift bins, directly testing the shape of the predicted curve.

\subsection{eROSITA}

The eROSITA X-ray survey~\citep{Merloni2012} provides X-ray hydrostatic mass
estimates for $\sim 100{,}000$ clusters. Cross-matching with Euclid and Rubin
weak lensing gives a uniform lensing-vs-X-ray comparison sample spanning
$0 < z < 1.5$.

\subsection{Combined test design}

The optimal test compares $M_\mathrm{lens}/M_\mathrm{dyn}$ in at least five
redshift bins spanning $0.1 < z < 2.0$ using a uniform mass estimation
methodology: (1)~select clusters with both weak lensing and spectroscopic
velocity dispersion measurements; (2)~apply consistent NFW mass models to both;
(3)~compute $M_\mathrm{lens}/M_\mathrm{dyn}$ in each bin with bootstrap errors;
(4)~fit to three models: constant (hydrostatic bias), IAM
(Equation~\ref{eq:ratio_explicit}, zero free parameters), and power law
$M_\mathrm{lens}/M_\mathrm{dyn} = 1 + A(1+z)^n$ (two free parameters); and
(5)~compare via BIC/AIC. The key discriminant is the functional form of the
$z$-dependence, not just its existence.

%-------------------------------------------------------------------
\section{Connection to IAM Cosmology}
%-------------------------------------------------------------------

The lensing-dynamics mass discrepancy is one manifestation of IAM's fundamental
prediction: the sector split between matter and photon observables. The same
$\mu < 1$, $\Sigma = 1$ signature produces the Hubble tension
($H_0^\mathrm{matter} = 72.26$~km\,s$^{-1}$\,Mpc$^{-1}$ vs.\ $H_0^\mathrm{photon}
= 67.16$~km\,s$^{-1}$\,Mpc$^{-1}$), the $\sigma_8$ suppression from $0.811$ to
$0.800$, and the $f\sigma_8(z)$ suppression at $z < 1$~\citep{Mahaffey2026a}.
All follow from $\beta_m = 0.1575$ and $E(a) = \exp(1-1/a)$ without additional
fitting. The lensing-dynamics test is particularly powerful because it is
measurable with current and near-future data, does not depend on CMB modeling,
and has a distinctive functional form in redshift.

%-------------------------------------------------------------------
\section{Discussion}
%-------------------------------------------------------------------

\subsection{Systematic concerns}

The primary systematic in any lensing-dynamics comparison is the calibration of
both mass estimates. Weak lensing masses depend on source redshift distributions,
shape measurement systematics, and projection effects. Dynamical masses depend on
assumptions about velocity anisotropy and dynamical equilibrium. X-ray masses
depend on temperature calibration and the hydrostatic assumption.

IAM's prediction applies to the \textit{ratio} $M_\mathrm{lens}/M_\mathrm{dyn}$,
which cancels many systematics that affect absolute mass estimates. The redshift
dependence provides additional discriminating power: systematics that evolve with
$z$ would need to conspire to mimic the specific functional form of
Equation~\ref{eq:ratio_explicit}, which is physically unmotivated for any known
systematic.

\subsection{Prior work on $\mu$--$\Sigma$ from cluster masses}

Several authors have constrained the $\mu/\Sigma$ ratio from cluster mass
comparisons~\citep{Terukina2014,Wilcox2015}. Pizzuti et al.~\citep{Pizzuti2017}
used cluster dynamics and lensing to constrain modified gravity models. These
analyses have generally focused on constant $\mu$, $\Sigma$ values rather than
redshift-dependent functions. IAM's specific prediction --- $\mu(a)$ declining
monotonically from $1$ at high $z$ to $0.864$ at $z = 0$, with $\Sigma = 1$
exactly --- has not been tested in this context.

\subsection{Falsifiability}

IAM's lensing-dynamics prediction is explicitly falsifiable. If
$M_\mathrm{lens}/M_\mathrm{dyn}$ shows no $z$-dependence, the hydrostatic bias
interpretation is preferred. If it increases with $z$, standard modified gravity
is preferred. If $M_\mathrm{lens}/M_\mathrm{dyn}$ at $z = 0$ differs from
$1.157$ by more than $3\sigma$ with sufficient sample size, the normalization is
challenged. If the $z$-dependence follows Equation~\ref{eq:ratio_explicit}, that
constitutes a direct measurement of $\mu(z)$ from cluster data. Whatever the
surveys find, the comparison to Equation~\ref{eq:ratio_explicit} is unambiguous.

%-------------------------------------------------------------------

\paragraph{Connection to the Dual-Sector Framework.}
The preservation of $\Sigma = 1$ derived here follows directly from the
dual-sector construction developed in the IAM cosmology paper.
Because the informational contribution renormalizes only the effective
matter-sector expansion rate, the photon sector and Weyl potential
relations remain identical to $\Lambda$CDM. This ensures internal
consistency across growth and lensing observables.

\section{Conclusion}
%-------------------------------------------------------------------

The Informational Actualization Model predicts a lensing-dynamics mass discrepancy
$M_\mathrm{lens}/M_\mathrm{dyn} = 1/\mu(z)$ of $15.7\%$ at $z = 0$, declining
monotonically to zero at high redshift following the specific, parameter-free
functional form of Equation~\ref{eq:ratio_explicit}. This prediction is consistent
with the 10--30\% lensing excess observed in multiple cluster surveys. The
redshift evolution provides a unique discriminant against hydrostatic bias
(flat in $z$), MOND (acceleration-dependent), and standard modified gravity
(opposite sign). Upcoming surveys from Euclid, Rubin, and eROSITA will provide
the statistical power to test the predicted redshift evolution at percent-level
precision, offering a direct measurement of $\mu(z)$ from cluster data.

%-------------------------------------------------------------------
\section*{Acknowledgments}
%-------------------------------------------------------------------

This work made use of NumPy~\citep{Harris2020}, SciPy~\citep{Virtanen2020},
and Matplotlib~\citep{Hunter2007}.

%-------------------------------------------------------------------
\begin{thebibliography}{99}

\bibitem{Amendola2008} Amendola, L., Kunz, M., \& Sapone, D.\ 2008,
JCAP, 04, 013

\bibitem{Biffi2016} Biffi, V., et al.\ 2016, ApJ, 827, 112

\bibitem{Euclid2022} Euclid Collaboration 2022, A\&A, 662, A112

\bibitem{Harris2020} Harris, C.\ R., et al.\ 2020, Nature, 585, 357

\bibitem{Hoekstra2015} Hoekstra, H., et al.\ 2015, MNRAS, 449, 685

\bibitem{Hunter2007} Hunter, J.\ D.\ 2007, Comput.\ Sci.\ Eng., 9, 90

\bibitem{Israel2014} Israel, H., et al.\ 2014, A\&A, 564, A129

\bibitem{Lau2009} Lau, E.\ T., Kravtsov, A.\ V., \& Nagai, D.\ 2009,
ApJ, 705, 1129

\bibitem{LSST2009} LSST Science Collaboration 2009, arXiv:0912.0201

\bibitem{Mahdavi2013} Mahdavi, A., et al.\ 2013, ApJ, 767, 116

\bibitem{Mahaffey2026a} Mahaffey, H.\ W.\ 2026a, ``Constraints on Late-Time
$f\sigma_8$ Suppression from $\mu < 1$, $\Sigma = 1$: Planck~2018 and
Large-Scale Structure,'' Universe, submitted (ID: universe-4189350)

\bibitem{Mahaffey2026b} Mahaffey, H.\ W.\ 2026b, ``Horizon Thermodynamics
and Gravitational Decoherence as the Origin of $\mu < 1$, $\Sigma = 1$,''
IAM Theory Paper, \url{https://github.com/hmahaffeyges/IAM-Validation/docs}

\bibitem{Merloni2012} Merloni, A., et al.\ 2012, arXiv:1209.3114

\bibitem{Nagai2007} Nagai, D., Kravtsov, A.\ V., \& Vikhlinin, A.\ 2007,
ApJ, 668, 1

\bibitem{Nelson2014} Nelson, K., et al.\ 2014, ApJ, 792, 25

\bibitem{PlanckClusters2016} Planck Collaboration 2016, A\&A, 594, A24

\bibitem{Planck2020lens} Planck Collaboration 2020, A\&A, 641, A8

\bibitem{Pizzuti2017} Pizzuti, L., et al.\ 2017, JCAP, 04, 023

\bibitem{Pogosian2016} Pogosian, L.\ \& Silvestri, A.\ 2016,
Phys.\ Rev.\ D, 94, 104014

\bibitem{Rasia2012} Rasia, E., et al.\ 2012, New J.\ Phys., 14, 055018

\bibitem{Smith2016} Smith, G.\ P., et al.\ 2016, MNRAS, 456, L74

\bibitem{Terukina2014} Terukina, A., et al.\ 2014, JCAP, 04, 013

\bibitem{Virtanen2020} Virtanen, P., et al.\ 2020, Nat.\ Methods, 17, 261

\bibitem{vonDerLinden2014} von der Linden, A., et al.\ 2014, MNRAS, 443, 1973

\bibitem{Wilcox2015} Wilcox, H., et al.\ 2015, MNRAS, 452, 1171

\end{thebibliography}

\end{document}
